\VerdierSemanticAnchor{front-editors-note}
\VerdierPhysicalPageAnchor{0008}
\chapter*{Editors' Note}
\addcontentsline{toc}{chapter}{Editors' Note}

\VerdierSemanticAnchor{front-editors-note-publication-history}
It was at the ceremony in tribute to J.-L. Verdier, organized by Université
Paris VII on 19 October 1989, that M. Artin, P. Cartier, and L. Illusie
conceived the idea of publishing J.-L. Verdier's thesis. Two years later I
joined them and, with the support of J. Le Potier and M. Broué, obtained funding
from Paris VII (URA 212) and the ÉNS (DMI) to key this text into \TeX{}. Once
the typing was completed, the result was carefully compared with Verdier's
typewritten original. A mathematical review still remained to be carried out,
since Verdier had never made the final corrections.

\VerdierSemanticAnchor{front-editors-note-chapter-i-modifications}
This careful reading was carried out, with the utmost care, by P. Cartier
(Chapter I), B. Keller (Chapter III), and, for the text as a whole, by
M. Zisman and myself. Apart from the many typographical errors that do not
merit enumeration, the principal modifications made are as follows. In
Chapter I, Proposition (3.1.5) asserted an equivalence. This equivalence was
replaced by an implication, since the converse (never used subsequently) is
false. Proposition (3.4.1) contained a sign error, giving rise to sign errors
in Chapter III ((1.3.4), (b), and (3.2.7.3)). This led to an adjustment of the
definition of the ``opposite triangulated category'' (Chapter II, 1.1.7), so
that the following assertion would remain true: \emph{the derived category of
the opposite category is the triangulated category opposite to the derived
category}.

\VerdierSemanticAnchor{front-editors-note-chapter-ii-modifications}
In Chapter II, the order of the propositions in §1.2 was changed. Indeed,
after writing the text Verdier had realized that the assertion that the direct
sum of two distinguished triangles is a distinguished triangle is a
consequence of the other axioms for triangulated categories. He had therefore
removed it from the list of axioms without making the resulting rearrangements.
In addition, paragraphs (2.3.6) and (2.3.7), which were not used later and
whose verification failed, were deleted. The definition of the notion of a
stationary spectral object (4.4.2) was slightly modified, the original
definition being too weak to establish the ensuing Theorem (4.4.3). It was
verified that this had no effect on what followed.

\VerdierPhysicalPageAnchor{0009}
\VerdierSemanticAnchor{front-editors-note-apparatus}
The index of notation and the terminological index were completed. By contrast,
the bibliography was left unchanged (except that the definitive references in
the ``Lecture Notes'' were given for the SGA volumes). In particular, the
absence of reference number 13 was not remedied.

\VerdierSemanticAnchor{front-editors-note-acknowledgements}
I would like to express my warm thanks to M. Broué, A. Bruguières, P. Cartier,
L. Gruson, L. Illusie, B. Keller, J. Le Potier, and M. Zisman, as well as to
Université Paris VII, École Normale Supérieure, and Institut des Hautes Études
Scientifiques, for the various ways in which they helped bring this project to
fruition. I also thank C. Defosse and C. Gourgues for patiently keying this work
into \textsc{AMS}-\TeX{}.

\VerdierSemanticAnchor{front-editors-note-memorial}
I would have liked to be able to extend these thanks to L. Doustaing, who died
prematurely in December 1994.

\VerdierSemanticAnchor{front-editors-note-signature}
\begin{flushright}
Paris, 14 June 1996\\
Georges Maltsiniotis
\end{flushright}
